%============================================================================
\section{Introducción}\label{sec:introduccion}
%============================================================================

El monitoreo de variables como temperatura, vibración o estado de máquina es
básico en la automatización industrial, y los enlaces inalámbricos que lo
soportan deben cumplir requisitos estrictos de confiabilidad~\cite{willig2005}
en entornos hostiles con interferencia de radiofrecuencia~\cite{gungor2009}. En
una planta, la conmutación de motores, contactores, variadores y cargas
inductivas produce interferencia impulsiva: eventos breves y de gran amplitud.
Los dispositivos con conmutación electromecánica, entre ellos los motores, se
han identificado como fuentes principales de este ruido~\cite{blackard1993}, que
es habitual en entornos industriales y degrada los sistemas
inalámbricos~\cite{chilo2009}; en canales sobre la red eléctrica incluye
componentes periódicas que hacen al canal variante en el tiempo~\cite{zimmermann2002}.

A diferencia del AWGN, estacionario y con potencia repartida en el tiempo, la
interferencia de origen humano es no gaussiana y degrada severamente a los
sistemas diseñados para ruido gaussiano~\cite{middleton1977}: concentra su
energía en pocas muestras de alta potencia~\cite{pighi2009}, y en portadora
única produce un piso de error que no desaparece al aumentar la energía de la
señal~\cite{ghosh1996}.

En este trabajo se estudió el escenario S1 del Proyecto Integrador de
Comunicaciones~II, representado en banda base: un enlace BPSK con pulso
rectangular, $R_s=\SI{32}{\kilo\sym\per\second}$, ocho muestras por símbolo y
$f_s=\SI{256}{\kilo\samples\per\second}$, al que se sumó un ruido impulsivo
periódico cuya amplitud fue la única variable, en una condición de referencia y
tres niveles de severidad.

El documento se organiza así: la Sección~\ref{sec:fundamento} presenta el
fundamento teórico; la Sección~\ref{sec:diseno}, el diseño del sistema; la
Sección~\ref{sec:metodologia}, la metodología de simulación; la
Sección~\ref{sec:resultados}, los resultados y su análisis; la
Sección~\ref{sec:hardware}, una prueba exploratoria en hardware SDR; y la
Sección~\ref{sec:conclusiones}, las conclusiones.
